% Advanced proof template
% Compile with: (XeLaTeX or LuaLaTeX recommended for unicode-math)
%   xelatex main.tex
% or (pdfLaTeX)
%   pdflatex main.tex
% If using minted, enable --shell-escape (optional)

\RequirePackage{iftex}
\documentclass[12pt,oneside]{book}

% -------------------------
% Encoding and fonts
% -------------------------
\usepackage[utf8]{inputenc} % safe for pdfLaTeX; ignored by Xe/Lua
\usepackage[T1]{fontenc}
\usepackage{lmodern}

% Engine-specific math font support
\ifXeTeX
  \usepackage{fontspec}
  \usepackage{unicode-math}
  \setmainfont{TeX Gyre Termes}
  \setmathfont{TeX Gyre Termes Math}
\else\ifLuaTeX
  \usepackage{fontspec}
  \usepackage{unicode-math}
  \setmainfont{TeX Gyre Termes}
  \setmathfont{TeX Gyre Termes Math}
\else
  % pdfLaTeX fallback
  \usepackage{amsfonts,amssymb}
  \usepackage{mathrsfs}
\fi\fi

% -------------------------
% Core math and theorem packages
% -------------------------
\usepackage{amsmath,amsthm,mathtools}
\usepackage{amssymb}
\usepackage{stmaryrd} % extra symbols
\usepackage{esint}    % extended integrals
\usepackage{bm}       % bold math symbols
\usepackage{nicefrac} % slanted fractions

% -------------------------
% Theorem styles
% -------------------------
\usepackage{thmtools}
\usepackage{thm-restate} % optional
\declaretheoremstyle[
  headfont=\bfseries,
  notefont=\normalfont,
  bodyfont=\normalfont,
  qed=\qedsymbol,
  headpunct={.}
]{mystyle}

\declaretheorem[style=mystyle,numberwithin=chapter]{theorem}
\declaretheorem[style=mystyle,sibling=theorem]{lemma}
\declaretheorem[style=mystyle,sibling=theorem]{proposition}
\declaretheorem[style=mystyle,sibling=theorem]{corollary}
\declaretheorem[style=mystyle,numbered=no]{remark}
\declaretheorem[style=mystyle,numbered=no]{example}
\declaretheorem[style=mystyle,numberwithin=chapter]{definition}
\declaretheorem[style=mystyle,numberwithin=chapter]{exercise}

% Proof environments: amsthm proof is available; we add variants
\usepackage{etoolbox}
\newenvironment{proofsketch}{\begin{proof}[Proof sketch]}{\end{proof}}
\newenvironment{proofidea}{\begin{proof}[Idea of proof]}{\end{proof}}

% Subproof environment for case analysis
\usepackage{chngcntr}
\newcounter{subproof}
\newenvironment{subproof}[1][]
  {\refstepcounter{subproof}\par\noindent\textbf{Subproof \thesubproof.}%
   \ifstrempty{#1}{}{\ \textit{(#1)}}\quad\ignorespaces}
  {\hfill$\square$\par}

% Boxed proof environment using mdframed
\usepackage[framemethod=TikZ]{mdframed}
\newmdenv[
  skipabove=6pt,
  skipbelow=6pt,
  leftmargin=0pt,
  rightmargin=0pt,
  innerleftmargin=8pt,
  innerrightmargin=8pt,
  backgroundcolor=gray!6,
  linecolor=black!30,
  linewidth=0.5pt
]{boxedproof}

% -------------------------
% Cross-referencing and TOC
% -------------------------
\usepackage{hyperref}
\hypersetup{
  colorlinks=true,
  linkcolor=blue,
  citecolor=blue,
  urlcolor=blue,
  pdftitle={Advanced Proofs Template},
  pdfauthor={}
}
\usepackage{cleveref}
\setcounter{tocdepth}{2} % include up to subsections in TOC
\setcounter{secnumdepth}{3}

% Optional mini-TOC per chapter
\usepackage{minitoc}
\dominitoc % initialize
\mtcsettitle{minitoc}{Contents of this chapter}

% -------------------------
% Graphics and plotting
% -------------------------
\usepackage{tikz}
\usetikzlibrary{arrows.meta,decorations.pathmorphing,positioning,graphs,graphs.standard}
\usepackage{pgfplots}
\pgfplotsset{compat=1.18}
\usepackage{tkz-graph} % graph utilities

% -------------------------
% Algorithms
% -------------------------
\usepackage{algorithm}
\usepackage{algpseudocode} % algorithmicx style
% alternative: algorithm2e
\usepackage[ruled,vlined]{algorithm2e}

% -------------------------
% Misc utilities
% -------------------------
\usepackage{enumitem}
\usepackage{microtype}
\usepackage{caption}
\usepackage{subcaption}
\usepackage{bookmark}
\usepackage{xcolor}
\usepackage{todonotes} % for notes during drafting

% -------------------------
% Custom macros
% -------------------------
\newcommand{\R}{\mathbb{R}}
\newcommand{\N}{\mathbb{N}}
\newcommand{\eps}{\varepsilon}
\newcommand{\abs}[1]{\left\lvert#1\right\rvert}
\newcommand{\norm}[1]{\left\lVert#1\right\rVert}
\newcommand{\set}[1]{\left\{#1\right\}}
\newcommand{\defeq}{\vcentcolon=}

% QED variants
\newcommand{\qedwhite}{\hfill\ensuremath{\square}}
\newcommand{\qedblack}{\hfill\ensuremath{\blacksquare}}

% -------------------------
% Document metadata
% -------------------------
\title{Advanced Proofs and Mathematical Writing}
\author{Author Name}
\date{\today}

% -------------------------
% Document start
% -------------------------
\begin{document}
\frontmatter
\maketitle
\tableofcontents
\listoffigures
\listoftables
\mainmatter

% Example chapter with minitoc
\chapter{Foundations}
\minitoc

\section{Notation and conventions}
We use standard notation: $\N$ for natural numbers, $\R$ for real numbers, etc.

\section{Sample theorem and proof types}

\begin{theorem}[Basic inequality]
Let $a,b\in\R$. Then $(a-b)^2\ge 0$ and hence $a^2+b^2\ge 2ab$.
\end{theorem}
\begin{proof}
Expand $(a-b)^2$ to get $a^2-2ab+b^2\ge 0$, rearrange to obtain the result.
\end{proof}

% -------------------------
% Proof by contrapositive
% -------------------------
\begin{theorem}
If $n^2$ is even then $n$ is even.
\end{theorem}
\begin{proof}[Proof by contrapositive]
We prove the contrapositive: if $n$ is odd then $n^2$ is odd. Write $n=2k+1$ and compute $n^2=4k^2+4k+1=2(2k^2+2k)+1$, which is odd.
\end{proof}

% -------------------------
% Proof by contradiction
% -------------------------
\begin{theorem}
There are infinitely many primes.
\end{theorem}
\begin{proof}[Proof by contradiction]
Assume finitely many primes $p_1,\dots,p_k$. Consider $N=\prod_{i=1}^k p_i +1$. Then $N$ is not divisible by any $p_i$, contradiction.
\end{proof}

% -------------------------
% Proof by induction
% -------------------------
\begin{theorem}
For all $n\in\N$, $\sum_{k=1}^n k = n(n+1)/2$.
\end{theorem}
\begin{proof}[Proof by induction]
Base $n=1$ holds. Assume true for $n$, then for $n+1$:


\[
\sum_{k=1}^{n+1} k = \frac{n(n+1)}{2} + (n+1) = \frac{(n+1)(n+2)}{2}.
\]


Thus holds for all $n$.
\end{proof}

% -------------------------
% Case analysis and subproofs
% -------------------------
\begin{theorem}
Let $x\in\R$. Then $x^2\ge 0$.
\end{theorem}
\begin{proof}
We consider cases.
\begin{itemize}
  \item If $x\ge 0$ then $x^2\ge 0$ trivially.
  \item If $x<0$ then let $y=-x>0$. Then $x^2=y^2\ge 0$.
\end{itemize}
\end{proof}

% -------------------------
% Constructive proof example
% -------------------------
\begin{theorem}
There exists an irrational number $x$ such that $x^x$ is rational.
\end{theorem}
\begin{proof}[Constructive proof]
Consider $x=\sqrt{2}$. If $\sqrt{2}^{\sqrt{2}}$ is rational we are done. Otherwise let $y=\sqrt{2}^{\sqrt{2}}$ which is irrational; then $y^{\sqrt{2}}=2$ is rational.
\end{proof}

% -------------------------
% Proof sketch and boxed proof
% -------------------------
\begin{theorem}
A sample boxed proof
\end{theorem}
\begin{boxedproof}
This is a proof placed inside a shaded box for emphasis. Use this style for long or important proofs.
\end{boxedproof}

% -------------------------
% Algorithms
% -------------------------
\chapter{Algorithms and Computation}
\section{Pseudocode example}
\begin{algorithm}
\caption{Euclid's algorithm for gcd}
\begin{algorithmic}[1]
\Procedure{GCD}{$a,b$}
  \While{$b\neq 0$}
    \State $t\gets b$
    \State $b\gets a \bmod b$
    \State $a\gets t$
  \EndWhile
  \State \textbf{return} $a$
\EndProcedure
\end{algorithmic}
\end{algorithm}

\section{Algorithm2e example}
\begin{algorithm}[H]
\SetAlgoLined
\KwIn{$a,b\in\mathbb{N}$}
\KwOut{$\gcd(a,b)$}
\While{$b\neq 0$}{
  $t\leftarrow b$\;
  $b\leftarrow a \bmod b$\;
  $a\leftarrow t$\;
}
\Return $a$\;
\caption{Euclid using algorithm2e}
\end{algorithm}

% -------------------------
% Graphs and plots
% -------------------------
\chapter{Graphs and Plots}
\section{TikZ graph example}
\begin{figure}[ht]
\centering
\begin{tikzpicture}[vertex/.style={circle,draw,inner sep=1.5pt}]
  \node[vertex] (A) at (0,0) {A};
  \node[vertex] (B) at (2,0) {B};
  \node[vertex] (C) at (1,1.5) {C};
  \draw[-{Stealth}] (A) -- (B);
  \draw[-{Stealth}] (B) -- (C);
  \draw[-{Stealth}] (C) -- (A);
\end{tikzpicture}
\caption{Directed triangle}
\end{figure}

\section{pgfplots example}
\begin{figure}[ht]
\centering
\begin{tikzpicture}
  \begin{axis}[
    axis lines=middle,
    xlabel=$x$, ylabel=$y$,
    domain=-2:2, samples=201,
    width=0.6\textwidth
  ]
    \addplot[blue,thick] {x^3 - x};
    \addplot[red,domain=-2:2] {sin(deg(x))};
    \legend{$x^3-x$,$\sin x$}
  \end{axis}
\end{tikzpicture}
\caption{Sample plots}
\end{figure}

% -------------------------
% Exotic math examples
% -------------------------
\chapter{Advanced Math Typesetting}
\section{Commutative diagrams and special symbols}
Use \verb|\mathscr| for script letters and \verb|\llbracket\rrbracket| from stmaryrd for interval notation.


\[
\mathscr{F}:\; 0 \to A \xrightarrow{f} B \xrightarrow{g} C \to 0
\]




\[
\llbracket a,b\rrbracket = \{x\in\Z: a\le x\le b\}
\]



% -------------------------
% Appendix and exercises
% -------------------------
\appendix
\chapter{Appendix}
\section{Useful identities}


\[
\int_{-\infty}^{\infty} e^{-x^2}\,dx = \sqrt{\pi}.
\]



\backmatter
\chapter{Exercises Solutions}
% ... solutions here

\end{document}
